% Generated by scripts/assemble.py; corrections belong in source/corrections.json.
\examhead{193}{云南大学 2024 年数学分析试卷}{182}
% p183-b22; source PDF page 183
\par\Needspace{5\baselineskip}
\blockanchor{p183-b22}
\majorquestion{一}填空题. 每题 5 分, 共 30 分.\par
% p183-b23; source PDF page 183
\blockanchor{p183-b23}
\question{1}极限 $\lim_{x\to 0}\left(\cot x - \frac{1}{x}\right) =$ \_\_\_\_.\par
% p183-b24; source PDF page 183
\blockanchor{p183-b24}
\question{2}设 $f(x) = \begin{cases} e^{ax}, & x \leq 0; \\ (x - b)^3, & x > 0. \end{cases}$ 在 $x = 0$ 处可导, 则 $a$ 和 $b$ 的值分别是 \_\_\_\_.\par
% p184-b00; source PDF page 184
\blockanchor{p184-b00}
\question{3}不定积分 $\int \frac{\mathrm{d}x}{2e^{-x} + e^x + 2} =$ \_\_\_\_.\par
% p184-b01; source PDF page 184
\blockanchor{p184-b01}
\question{4}曲面 $z = x + 2y + \ln (1 + x^2 + y^2)$ 在 $(0,0,0)$ 的切平面方程是\par
% p184-b02; source PDF page 184
\blockanchor{p184-b02}
\question{5}设 $f(x, y) = \int_{0}^{xy} e^{xt^2} \, \mathrm{d}t$，则 $\left. \frac{\partial^2 f}{\partial x \partial y} \right|_{(1,1)} =$ \_\_\_\_.\par
% p184-b03; source PDF page 184
\blockanchor{p184-b03}
\question{6}曲线 $y = \int_0^x \tan t \, \mathrm{d}t, 0 \leq x \leq \frac{\pi}{4}$ 的弧长 $s =$ \_\_\_\_.\par
% p184-b04; source PDF page 184
\par\Needspace{5\baselineskip}
\blockanchor{p184-b04}
\majorquestion{二}(15 分) 按定义证明: $\lim_{x \to 1} \frac{x(x - 1)}{x^2 - 1} = \frac{1}{2}$.\par
% p184-b05; source PDF page 184
\par\Needspace{5\baselineskip}
\blockanchor{p184-b05}
\majorquestion{三}(15 分) 判别级数 $\sum_{n=1}^{\infty} \frac{a^n n!}{n^n} (a > 0)$ 的敛散性.\par
% p184-b06; source PDF page 184
\par\Needspace{5\baselineskip}
\blockanchor{p184-b06}
\majorquestion{四}(15 分) 设函数 $f(x)$ 和 $g(x)$ 在区间 $[a, b]$ 上二阶可导, 并且 $g''(x) \neq 0$, $f(a) = f(b) = g(a) = g(b) = 0$. 证明:\par
% p184-b07; source PDF page 184
\blockanchor{p184-b07}
(1) 在开区间 $(a, b)$ 内, $g(x) \neq 0$.\par
% p184-b08; source PDF page 184
\blockanchor{p184-b08}
(2) 在开区间 $(a, b)$ 内, 至少存在一点 $\xi$, 使得 $\frac{f(\xi)}{g(\xi)} = \frac{f''(\xi)}{g''(\xi)}$.\par
% p184-b09; source PDF page 184
\par\Needspace{5\baselineskip}
\par\Needspace{6\baselineskip}
\blockanchor{p184-b09}
\majorquestion{五}(15 分) 已知函数 $z = ue^{ax + by}$, 其中 $u = u(x, y)$ 具有二阶偏导数, 且 $\frac{\partial^2 u}{\partial x \partial y} = 0$. 确定常数 $a, b$, 使得函数 $z = z(x, y)$ 满足方程\par
\nopagebreak[4]
% p184-b10; source PDF page 184
\blockanchor{p184-b10}
\displayformula{\frac {\partial^ {2} z}{\partial x \partial y} + 2 \frac {\partial z}{\partial x} + 3 \frac {\partial z}{\partial y} + 6 z = 0.}
% p184-b11; source PDF page 184
\par\Needspace{5\baselineskip}
\par\Needspace{6\baselineskip}
\blockanchor{p184-b11}
\majorquestion{六}(15 分) 设 $f(x)$ 在 $t \neq 0$ 时一阶连续可导, 且 $f(1) = 0$. 求函数 $f(x^2 + y^2)$, 使得曲线积分\par
\nopagebreak[4]
% p184-b12; source PDF page 184
\blockanchor{p184-b12}
\displayformula{\int_ {L} y (2 - f (x ^ {2} + y ^ {2})) \mathrm{d} x + x f (x ^ {2} + y ^ {2}) \mathrm{d} y}
% p184-b13; source PDF page 184
\blockanchor{p184-b13}
与路径无关, 其中 L 为不经过原点的光滑曲线.\par
% p184-b14; source PDF page 184
\par\Needspace{5\baselineskip}
\blockanchor{p184-b14}
\majorquestion{七}(15 分) 求级数 $\sum_{n=1}^{\infty} n(n+2)x^{n-1}$ 的和函数.\par
% p184-b15; source PDF page 184
\par\Needspace{5\baselineskip}
\blockanchor{p184-b15}
\majorquestion{八}(15 分) 计算极限 $\lim_{a \to +\infty} \iint_{D} \min\{x, y\} e^{-(x^2 + y^2)} \, \mathrm{d}x \, \mathrm{d}y$, 其中 $D: [-a, a] \times [-a, a]$.\par
% p184-b16; source PDF page 184
\par\Needspace{5\baselineskip}
\blockanchor{p184-b16}
\majorquestion{九}(15 分) 设 $f(x)$ 在 $[a, +\infty)$ 上一致连续, $g(x)$ 在 $[a, +\infty)$ 上连续, 且 $\lim_{x \to +\infty} [f(x) - g(x)] = 0$. 证明: $g(x)$ 在 $[a, +\infty)$ 上一致连续.\par
